% Luc Destombe --- licence CC BY-NC-SA 4.0
% https://creativecommons.org/licenses/by-nc-sa/4.0/deed.fr
% Fichier autonome : compiler avec pdflatex, deux fois (signets, nombre de pages).
\documentclass[10pt,a4paper,twocolumn]{article}
\makeatletter
\newif\iflycee@unecolonne
\newif\iflycee@palatino
\lycee@palatinotrue

\RequirePackage[utf8]{inputenc}
\RequirePackage[T1]{fontenc}
\RequirePackage[french]{babel}
\RequirePackage{etoolbox}

\newcommand{\ShorthandsOff}{\@ifundefined{shorthandoff}{}{\shorthandoff{;:!?}}}
\newcommand{\ShorthandsOn}{\@ifundefined{shorthandon}{}{\shorthandon{;:!?}}}
\ShorthandsOff
\AtBeginDocument{\ShorthandsOn}
\AtBeginEnvironment{tikzpicture}{\ShorthandsOff}

\RequirePackage{geometry}
\RequirePackage{fancyhdr}
\RequirePackage{lastpage}
\RequirePackage{multicol}
\RequirePackage{array}
\RequirePackage{enumitem}
\RequirePackage{xcolor}
\RequirePackage{needspace}

\RequirePackage{amsmath,amssymb}
\RequirePackage{mathpazo}
\DeclareMathSizes{10.95}{10.95}{8}{5.5}
\begingroup\catcode`\_=\active \catcode`\^=\active
\gdef_#1{\sb{\scriptscriptstyle #1}}
\gdef^#1{\sp{\scriptscriptstyle\lycee@moinsepais #1}}
\endgroup
\newcommand\lycee@moins{\mathbin{%
  \setbox\z@\hbox{$\m@th\scriptscriptstyle\lycee@moinsorig$}%
  \dimen@=.55\wd\z@ \advance\dimen@-.5pt
  \hbox to.75\wd\z@{\hss$\m@th\scriptscriptstyle\vcenter{\hbox{%
    \tikz\draw[line width=.5pt,line cap=round](0,0)--(\the\dimen@,0);}}$\hss}}}
\begingroup\lccode`\~=`\- \lowercase{\endgroup\let~}\lycee@moins
\newcommand\lycee@moinsepais{\mathcode`\-="8000 }
\AtBeginDocument{\mathchardef\lycee@moinsorig=\mathcode`\-\relax}
\AtBeginDocument{\catcode`\_=12 \mathcode`\_="8000
  \catcode`\^=12 \mathcode`\^="8000 }
\newcommand\lycee@exposants{%
  \fontdimen13\textfont2=.48\fontdimen6\textfont2
  \fontdimen14\textfont2=.48\fontdimen6\textfont2
  \fontdimen15\textfont2=.48\fontdimen6\textfont2
  \fontdimen13\scriptfont2=.48\fontdimen6\scriptfont2
  \fontdimen14\scriptfont2=.48\fontdimen6\scriptfont2
  \fontdimen15\scriptfont2=.48\fontdimen6\scriptfont2}
\every@math@size\expandafter{\the\every@math@size\lycee@exposants}
\newlength{\retraitcalculs}
\setlength{\retraitcalculs}{2em}

\RequirePackage{tikz}
\usetikzlibrary{arrows.meta,calc,babel,shapes.misc}
\RequirePackage[most]{tcolorbox}
\tcbuselibrary{skins,breakable}

\definecolor{coulDef}{HTML}{1F4E79}
\definecolor{coulProp}{HTML}{7030A0}
\definecolor{coulRem}{HTML}{C55A11}
\definecolor{coulEx}{HTML}{2E7D32}
\definecolor{coulExo}{HTML}{B03A2E}
\definecolor{coulPart}{HTML}{1F4E79}
\definecolor{coulMeth}{HTML}{1B6E4F}
\definecolor{coulCourbe}{HTML}{0B5ED7}
\definecolor{coulCourbeB}{HTML}{C00000}
\definecolor{coulCourbeC}{HTML}{2E7D32}
\definecolor{coulGrille}{HTML}{8C8C8C}
\definecolor{coulRep}{HTML}{1B6E4F}
\definecolor{coulBar}{HTML}{2E7D32}

\newcommand{\R}{\mathbb{R}}
\newcommand{\Rs}{\mathbb{R}^{*}}
\renewcommand{\le}{\leqslant}
\renewcommand{\ge}{\geqslant}
\renewcommand{\approx}{\simeq}
\newcommand{\vect}[1]{\overrightarrow{#1}}
\newcommand{\norme}[1]{\left\lVert #1 \right\rVert}
\newcommand{\intervalle}[2]{\left[#1\,;#2\right]}
\RequirePackage{cancel}
\renewcommand{\CancelColor}{\color{coulExo}}
\newcommand{\fois}[1]{\times{\color{coulExo}#1}}
\newcommand{\barre}[1]{\cancel{#1}}

\tikzset{grille/.style={coulGrille,line width=0.4pt}}
\newcommand{\quadrillage}[4]{\draw[grille,step=1] (#1,#2) grid (#3,#4);}
\tikzset{
  croix/.style={cross out,draw,line width=0.6pt,inner sep=0pt,minimum size=4pt},
  croixdroite/.style={croix,rotate=45,line width=0.8pt},
}
\newcommand{\pt}[3]{\path (#1) node[croix]{} node[#2,font=\small,inner sep=2.5pt]{$#3$};}

\newcommand{\lycee@repere}[4]{%
  \draw[grille,step=1] (#1,#2) grid (#3,#4);
  \draw[-{Stealth[length=2mm]},thick] (#1,0)--({#3+0.4},0) node[below left=-1pt]{$x$};
  \draw[-{Stealth[length=2mm]},thick] (0,#2)--(0,{#4+0.4}) node[below left=-1pt]{$y$};
  \node[below left,font=\scriptsize,inner sep=1.5pt] at (0,0) {$0$};}
\newcommand{\lycee@gradx}[1]{\foreach \k/\l in {#1}{%
  \draw (\k,0.07)--(\k,-0.07) node[below,font=\scriptsize,inner sep=1.5pt]{$\l$};}}
\newcommand{\lycee@grady}[1]{\foreach \k/\l in {#1}{%
  \draw (0.07,\k)--(-0.07,\k) node[left,font=\scriptsize,inner sep=1.5pt]{$\l$};}}
\newcommand{\lycee@intff}[2]{\left[#1\,;#2\right]}
\newcommand{\lycee@intfo}[2]{\left[#1\,;#2\right[}
\newcommand{\lycee@intof}[2]{\left]#1\,;#2\right]}
\newcommand{\lycee@intoo}[2]{\left]#1\,;#2\right[}
\newcommand{\lycee@ens}[1]{\left\{#1\right\}}
\AtBeginDocument{%
  \providecommand{\repere}{\lycee@repere}%
  \providecommand{\gradx}{\lycee@gradx}%
  \providecommand{\grady}{\lycee@grady}%
  \providecommand{\Cf}{\mathcal{C}\sb{\scriptscriptstyle f}}%
  \providecommand{\Cg}{\mathcal{C}\sb{\scriptscriptstyle g}}%
  \providecommand{\intff}{\lycee@intff}%
  \providecommand{\intfo}{\lycee@intfo}%
  \providecommand{\intof}{\lycee@intof}%
  \providecommand{\intoo}{\lycee@intoo}%
  \providecommand{\ens}{\lycee@ens}}

\newlist{questions}{enumerate}{3}
\setlist[questions,1]{label=\arabic*\textdegree),leftmargin=*,itemsep=2pt,topsep=3pt}
\setlist[questions,2]{label=\alph*),leftmargin=*,itemsep=1pt,topsep=2pt}
\setlist[questions,3]{label=\roman*),leftmargin=*,itemsep=1pt,topsep=1pt}
\setlist[itemize]{label=\textbullet,leftmargin=*,itemsep=2pt,topsep=3pt}

\newcommand{\besoinplace}[1]{\par\penalty\@M\begingroup
  \dimen@=#1\relax\dimen@ii=\pagegoal\advance\dimen@ii-\pagetotal
  \ifdim\dimen@>\dimen@ii\ifdim\dimen@ii>\z@\vfil\fi\break\fi\endgroup}

\newcommand{\lycee@pied}{\thepage/\pageref{LastPage}}
\newcommand{\entete}[2]{%
  \pagestyle{fancy}\fancyhf{}%
  \fancyhead[L]{\footnotesize\color{coulPart}\textbf{#1}}%
  \fancyhead[R]{\footnotesize\color{coulPart}\textbf{#2}}%
  \fancyfoot[C]{\footnotesize\lycee@pied}%
  \renewcommand{\headrulewidth}{0.4pt}%
  \renewcommand{\headrule}{\hbox to\headwidth{%
    \color{coulPart!40}\leaders\hrule height \headrulewidth\hfill}}}

\RequirePackage{ccicons}
\newcommand{\lycee@licence}{{\scriptsize\color{gray}%
  \copyright\ Luc Destombe --- \ccbyncsa\ CC BY-NC-SA 4.0}}
\newif\iflycee@licence \lycee@licencetrue
\newcommand{\sanslicence}{\lycee@licencefalse}
\AtBeginDocument{\iflycee@licence\AddToHookNext{shipout/foreground}{%
  \put(0,0){\rlap{\kern\dimexpr1in+\hoffset+\oddsidemargin+\textwidth\relax
    \llap{\raisebox{-\dimexpr1in+\voffset+\topmargin+\headheight+\headsep
      +\textheight+\footskip\relax}{\lycee@licence}}}}}\fi}

\newcommand{\formulesserrees}{%
  \setlength{\abovedisplayskip}{3pt}\setlength{\belowdisplayskip}{3pt}%
  \setlength{\abovedisplayshortskip}{2pt}\setlength{\belowdisplayshortskip}{3pt}}

\setlength{\parindent}{0pt}

\RequirePackage[implicit=false,hidelinks,pdfencoding=auto,psdextra,bookmarksopen,
  bookmarksopenlevel=1,pdfcreator={},pdfproducer={}]{hyperref}
\RequirePackage{bookmark}
\hypersetup{pdfauthor={Luc Destombe},pdfkeywords={CC BY-NC-SA 4.0}}
\newcounter{lycee@signet}
\newcommand{\signet}[3][]{\stepcounter{lycee@signet}%
  \edef\lycee@dest{\if\relax\detokenize{#1}\relax
    signet.\arabic{lycee@signet}\else#1\fi}%
  \hypertarget{\lycee@dest}{}\bookmark[level=#2,dest=\lycee@dest]{#3}}
\pdfstringdefDisableCommands{%
  \def\R{R}\def\vect#1{#1}\def\Cf{Cf}\def\Cg{Cg}\def\fois#1{×#1}%
  \def\barre#1{#1}\def\intervalle#1#2{[#1;#2]}\def\intff#1#2{[#1;#2]}%
  \def\textsuperscript#1{#1}\def\dfrac#1#2{#1/#2}\def\frac#1#2{#1/#2}%
  \def\vec#1{#1⃗}}%
\RequirePackage[official]{eurosym}

\tcbset{exercice corrige/.style={
  enhanced,breakable,before upper=\raggedright,
  colback=white,colframe=coulExo,
  boxrule=0.8pt,arc=2pt,
  left=6pt,right=6pt,top=8pt,bottom=5pt,
  before skip=9pt,after skip=4pt,
  coltitle=white,fonttitle=\bfseries\small,
  attach boxed title to top left={xshift=8pt,yshift=-\tcboxedtitleheight/2},
  boxed title style={colback=coulExo,arc=2pt,boxrule=0pt}}}
\newcounter{exo}
\newtcolorbox{exercice}[1][]{exercice corrige,
  phantom={\signet[exercice-\arabic{exo}]{2}{Exercice \arabic{exo}}},
  title={Exercice \theexo\ifx\relax#1\relax\else{} --- #1\fi}}
\newenvironment{exo}[1][\relax]{\refstepcounter{exo}\begin{exercice}[#1]}{\end{exercice}}

\RequirePackage{cuted}
\geometry{top=1.6cm,bottom=1cm,left=1cm,right=1cm,headheight=14pt,footskip=0.6cm}

\newcommand{\entetecorrige}[3]{%
  \begin{strip}
  \begin{center}
    \begin{tcolorbox}[enhanced,width=0.92\linewidth,colback=white,colframe=coulPart,
        boxrule=1.6pt,arc=3pt,halign=center,top=5pt,bottom=5pt]
      {\LARGE\bfseries\color{coulPart} #1}\\[3pt]
      {\normalsize #2}
    \end{tcolorbox}
  \end{center}
  \if\relax\detokenize{#3}\relax\else

  \vspace{2pt}
  \noindent
  \begin{tcolorbox}[enhanced,colback=white,colframe=black!35,boxrule=0.5pt,arc=2pt,
      left=7pt,right=7pt,top=4pt,bottom=4pt]
  {\small #3}
  \end{tcolorbox}
  \vspace{6pt}
  \fi
  \end{strip}}

\newcounter{partie}
\newcommand{\partiebandeau}[1]{%
  \besoinplace{8\baselineskip}\vspace{6pt}%
  \begin{tcolorbox}[enhanced,colback=coulPart!8,colframe=coulPart,boxrule=0pt,
      leftrule=3pt,arc=0pt,left=5pt,right=4pt,top=3pt,bottom=3pt,
      before skip=4pt,after skip=2pt]
    \bfseries\color{coulPart}#1\signet{1}{#1}
  \end{tcolorbox}}
\newcommand{\partie}[1]{\stepcounter{partie}\partiebandeau{\Roman{partie}) #1}}
\newcommand{\partiesansnum}[1]{\partiebandeau{#1}}

\newcommand{\res}[1]{%
  \tcbox[on line,colback=coulPart!7,colframe=coulPart,boxrule=0.5pt,arc=1pt,
    boxsep=0pt,left=3pt,right=3pt,top=1pt,bottom=2pt]{$#1$}}
\newcommand{\calc}[1]{\par\smallskip\noindent$\displaystyle #1$\par}
\newenvironment{calculs}{\par\smallskip\noindent$\displaystyle\begin{aligned}[t]}%
  {\end{aligned}$\par\smallskip}
\newcommand{\rem}[1]{\par\smallskip{\small\itshape\color{coulPart}#1}\par}
\newcommand{\deux}[2]{\par\noindent\makebox[0.5\linewidth][l]{#1}#2}

\setlist[questions,1]{label=\arabic*\textdegree),leftmargin=*,itemsep=2pt,topsep=2pt}
\setlist[questions,2]{label=\alph*),leftmargin=*,itemsep=2pt,topsep=1pt}
\newlist{reponses}{enumerate}{1}
\setlist[reponses]{label=\alph*),leftmargin=*,itemsep=2pt,topsep=2pt}

\setlength{\parskip}{1pt}
\setlength{\columnsep}{0.6cm}
\raggedbottom
\formulesserrees
\AtBeginDocument{\formulesserrees}

\makeatother
\entete{Chapitre 3 --- Vecteurs --- Corrigé des exercices}{Seconde}

\let\deuxpaquet\deux
\renewcommand{\deux}[2]{\deuxpaquet{#1}{#2}\par}
\tikzset{
  vecteur/.style={-{Stealth[length=2.2mm]},line width=0.9pt,coulCourbe},
  vecteurB/.style={-{Stealth[length=2.2mm]},line width=0.9pt,coulCourbeB},
  vecteurC/.style={-{Stealth[length=2.2mm]},line width=0.9pt,coulCourbeC},
  aide/.style={dashed,black!50,line width=0.5pt},
  fig/.style={line width=0.7pt,black!75},
}
\clubpenalty=10000
\widowpenalty=10000

\begin{document}

\entetecorrige{CORRIGÉ --- VECTEURS}
  {Chapitre 3 : fiche d'exercices --- Seconde}
  {Les résultats sont encadrés ; les remarques en bleu signalent les erreurs
   fréquentes et les méthodes à retenir. Sur un quadrillage, on lit chaque vecteur comme
   un déplacement : tant de carreaux à droite ou à gauche, tant vers le haut ou le bas.}

\partie{Caractéristiques d'un vecteur}

\begin{exo}[Égaux ou non]
\deux{a) $\vect{KL}\res{=}\vect{NP}$}{b) $\vect{KN}\res{=}\vect{LP}$}
\deux{c) $\vect{KP}\res{=}\vect{LQ}$}{d) $\vect{NL}\res{\neq}\vect{KN}$}
\deux{e) $\vect{KM}\res{=}\vect{NQ}$}{f) $\vect{PQ}\res{\neq}\vect{PN}$}
\rem{f) Même direction, même norme, mais sens contraires : $\vect{PN}=-\vect{PQ}$.}
\end{exo}

\begin{exo}[Lire sur la figure]
1) $\vect{KN}=\res{\vect{LP}}=\vect{MQ}$ ; \quad $\vect{LM}=\res{\vect{KL}}=\vect{NP}=\vect{PQ}$.\par
2) Opposés à $\vect{KN}$ : $\res{\vect{NK}}$, $\vect{PL}$, $\vect{QM}$ ; à $\vect{KL}$ :
$\res{\vect{LK}}$, $\vect{PN}$, $\vect{QP}$\dots\par
3) Par exemple : $\vect{KL}$ et $\vect{KM}$ ; \; $\vect{KL}$ et $\vect{PN}$ ; \;
$\vect{KM}$ et $\vect{LK}$.
\end{exo}

\begin{exo}[Construire un représentant]
\begin{center}
\begin{tikzpicture}[scale=0.55]
  \draw[vecteur] (0,3)--(3,3) node[midway,above,font=\small]{$\vec{u}$};
  \draw[vecteurB] (0,2)--(6,2) node[midway,above,font=\small]{1)};
  \draw[vecteurB] (3,1)--(0,1) node[midway,above,font=\small]{2)};
  \draw[vecteurB] (1.5,0)--(0,0) node[midway,above,font=\small]{3)};
  \draw[vecteurB] (4,0)--(5.8,-2.4) node[midway,right,font=\small]{4)};
\end{tikzpicture}
\end{center}
\rem{1) et 3) : on mesure la longueur demandée sur la même direction. 4) Toute flèche
de $3$~cm dans une autre direction convient.}
\end{exo}

\begin{exo}[Vrai ou faux]
1) \res{\text{Vrai}} : $EFGH$ est un parallélogramme, donc $\vect{GF}=\vect{HE}$, et la
translation de vecteur $\vect{HE}$ envoie $H$ sur $E$.\par
2) \res{\text{Vrai}} : $HGKL$ est un parallélogramme, donc $(GK)\parallel(HL)$.\par
3) \res{\text{Faux}} : $\vect{GK}=\vect{HL}$, donc $\vect{KG}=-\vect{HL}$ : les vecteurs
sont opposés.\par
4) \res{\text{Vrai}} : $\vect{EF}=\vect{HG}$ et $\vect{HG}=\vect{LK}$.
\end{exo}

\begin{exo}[Dans un parallélogramme]
1) $\vect{RS}=\res{\vect{UT}}$ ; \; $\vect{TU}=\res{\vect{SR}}$ ; \; $\vect{RU}=\res{\vect{ST}}$.\par
2) $\vect{RO}=\res{\vect{OT}}$ ; \; $\vect{UO}=\res{\vect{OS}}$ ; \; $\vect{SO}=\res{\vect{OU}}$.
\rem{Les diagonales $[RT]$ et $[SU]$ se coupent en leur milieu $O$.}
\end{exo}

\partie{Translation}

\begin{exo}[Représentants d'origine donnée]
\begin{center}
\begin{tikzpicture}[scale=0.5]
  \draw[fig] (0,0)--(4,0)--(1,3)--cycle;
  \pt{0,0}{below left}{D} \pt{4,0}{below right}{E} \pt{1,3}{above}{F}
  \pt{6,2}{left}{P}
  \draw[vecteur] (6,2)--(10,2) node[midway,above,font=\scriptsize]{$\vect{DE}$};
  \draw[vecteurB] (6,2)--(3,5) node[midway,left,font=\scriptsize]{$\vect{EF}$};
  \draw[vecteurC] (6,2)--(5,-1) node[midway,left,font=\scriptsize]{$\vect{FD}$};
\end{tikzpicture}
\end{center}
\rem{Chaque flèche part de $P$, avec la direction, le sens et la longueur du côté.}
\end{exo}

\begin{exo}[Quatre points à placer]
\begin{center}
\begin{tikzpicture}[scale=0.42]
  \quadrillage{-3.4}{-3.4}{5.4}{3.4}
  \draw[fig] (0,0)--(4,0)--(1,3)--cycle;
  \pt{0,0}{below left}{A} \pt{4,0}{below right}{B} \pt{1,3}{above}{C}
  \draw[vecteur] (4,0)--(5,3); \pt{5,3}{above right}{M}
  \draw[vecteur] (3,-3)--(4,0); \pt{3,-3}{below}{N}
  \draw[vecteurB] (-3,3)--(1,3); \pt{-3,3}{above}{P}
  \draw[vecteurC] (0,0)--(3,-3);
\end{tikzpicture}
\end{center}
\rem{$Q$ est confondu avec $N$ : $\vect{AQ}=\vect{CB}$ et $\vect{NB}=\vect{AC}$ donnent le
même point. $ABMC$, $ANBC$ et $APCB$ sont des parallélogrammes.}
\end{exo}

\begin{exo}[Origine ou extrémité]
\begin{center}
\begin{tikzpicture}[scale=0.5]
  \draw[vecteur] (0,3)--(2,4) node[midway,above,font=\small]{$\vec{u}$};
  \pt{0,0}{below left}{E} \draw[vecteurB] (0,0)--(2,1);
  \pt{6,3}{above right}{F} \draw[vecteurB] (4,2)--(6,3);
  \pt{7,0}{below right}{G} \draw[vecteurB] (5,-1)--(7,0);
\end{tikzpicture}
\end{center}
\rem{Pour un représentant d'extrémité $F$, on part de $F$ en reculant du déplacement de
$\vec{u}$, puis on trace la flèche vers $F$.}
\end{exo}

\begin{exo}[Image d'un triangle]
On reporte le vecteur $\vect{BC}$ depuis $A$, $B$ et $C$ ; l'image de $B$ est $C$.
Les triangles $ABC$ et $A'B'C'$ sont \res{\text{superposables}} : une translation
conserve les longueurs ($A'B'=5$~cm, $B'C'=4$~cm, $A'C'=3$~cm).
\end{exo}

\begin{exo}[Neuf points]
1) Égaux à $\vect{AB}$ : \res{\vect{BC}\,;\vect{DE}\,;\vect{EF}\,;\vect{GH}\,;\vect{HI}}\par
2) Égaux à $\vect{AE}$ : \res{\vect{BF}\,;\vect{DH}\,;\vect{EI}}\par
3) Par exemple \res{\vect{BA}} et \res{\vect{CB}}.
\end{exo}

\partie{Parallélogramme et égalité de vecteurs}

\begin{exo}[Deux parallélogrammes]
1) $EFGH$ est un parallélogramme, donc $\vect{EF}=\res{\vect{HG}}$ ; $EFIJ$ en est un,
donc $\vect{EF}=\res{\vect{JI}}$.\par
2) Ainsi $\vect{HG}=\vect{JI}$ ; d'après le théorème 1, \res{HGIJ} est un
parallélogramme.
\rem{$\vect{AB}=\vect{CD}$ donne $ABDC$ : ici $\vect{HG}=\vect{JI}$ donne $HGIJ$.}
\end{exo}

\begin{exo}[Autour d'un triangle]
\begin{center}
\begin{tikzpicture}[scale=0.42]
  \quadrillage{-3.4}{-3.4}{5.4}{3.4}
  \draw[fig] (0,0)--(4,0)--(1,3)--cycle;
  \pt{0,0}{below left}{D} \pt{4,0}{below right}{E} \pt{1,3}{above}{F}
  \draw[vecteur] (1,3)--(5,3); \pt{5,3}{above right}{M}
  \draw[vecteur] (-3,3)--(1,3); \pt{-3,3}{above}{N}
  \draw[vecteurB] (0,0)--(3,-3); \pt{3,-3}{below}{K}
\end{tikzpicture}
\end{center}
3) C'est \res{M} : $\vect{ME}=\vect{FD}$ car $FDEM$ est un parallélogramme.\par
4) Par exemple $\vect{NF}=\vect{FM}$ ; $\vect{DK}=\vect{FE}$ ; $\vect{DF}=\vect{EM}$.
\end{exo}

\begin{exo}[Un milieu à démontrer]
$RSTU$ est un parallélogramme, donc $\vect{RS}=\vect{UT}$. $RSUV$ en est un, donc
$\vect{RS}=\vect{VU}$. Ainsi $\vect{VU}=\vect{UT}$ : \res{U \text{ est le milieu de } [TV]}.
\end{exo}

\begin{exo}[Vrai ou faux]
1) \res{\text{Faux}} : c'est $\vect{CF}=\vect{DE}$ (le vecteur va de $C$ vers son
image).\par
2) \res{\text{Vrai}} : $\vect{KL}=\vect{NM}$ donne le parallélogramme $KLMN$
(théorème 1).\par
3) \res{\text{Vrai}} : $\vect{AD}=\vect{BC}$, donc la translation envoie $B$ sur $C$.
\end{exo}

\partie{Somme de deux vecteurs}

\begin{exo}[Trois points à construire]
\begin{minipage}[t]{0.42\linewidth}\raggedright
$\vect{BC}+\vect{CA}=\vect{BA}$ : \res{D=A}.\par
On reporte $\vect{BC}$ à la suite de $\vect{BA}$, depuis $A$ : $E$ ; $BAEC$ est un
parallélogramme.\par
$\vect{CF}=\vect{CB}+\vect{AB}$ : on met bout à bout.
\end{minipage}\hfill
\begin{minipage}[t]{0.55\linewidth}
\vspace{0pt}\centering
\begin{tikzpicture}[scale=0.38]
  \quadrillage{-3.4}{-3.4}{8.4}{3.4}
  \draw[fig] (0,0)--(4,0)--(1,3)--cycle;
  \pt{0,0}{below left}{A} \pt{4,0}{below right}{B} \pt{1,3}{above right}{C}
  \draw[aide] (0,0)--(-3,3)--(1,3);
  \pt{-3,3}{above}{E}
  \draw[vecteurB] (1,3)--(4,0); \draw[vecteurB] (4,0)--(8,0);
  \pt{8,0}{below}{F}
\end{tikzpicture}
\end{minipage}
\end{exo}

\begin{exo}[Somme et différence sur un quadrillage]
\begin{center}
\begin{tikzpicture}[scale=0.42]
  \quadrillage{-0.4}{-0.4}{5.4}{4.4}
  \pt{1,1}{below left}{A}
  \draw[vecteur] (1,1)--(4,2); \draw[vecteurB] (1,1)--(2,3);
  \draw[aide] (4,2)--(5,4); \draw[vecteurC] (1,1)--(5,4); \pt{5,4}{above}{M}
  \draw[aide] (4,2)--(3,0); \draw[vecteurC] (1,1)--(3,0); \pt{3,0}{below}{N}
\end{tikzpicture}
\hspace{4pt}
\begin{tikzpicture}[scale=0.42]
  \quadrillage{-0.4}{-0.4}{5.4}{4.4}
  \pt{1,2}{left}{A}
  \draw[vecteur] (1,2)--(4,2); \draw[vecteurB] (1,2)--(1,0);
  \draw[aide] (4,2)--(4,0); \draw[vecteurC] (1,2)--(4,0); \pt{4,0}{below}{M}
  \draw[aide] (4,2)--(4,4); \draw[vecteurC] (1,2)--(4,4); \pt{4,4}{above}{N}
\end{tikzpicture}
\end{center}
\rem{Gauche : $\vec{u}+\vec{v}$ fait $4$ à droite et $3$ vers le haut ;
$\vec{u}-\vec{v}$, $2$ à droite et $1$ vers le bas.}
\end{exo}

\begin{exo}[Origine imposée]
\begin{center}
\begin{tikzpicture}[scale=0.42]
  \quadrillage{-0.4}{-0.4}{5.4}{4.4}
  \pt{1,1}{below left}{T}
  \draw[vecteur] (0,2)--(1,4); \draw[vecteurB] (3,4)--(5,3);
  \draw[aide] (1,1)--(2,3)--(4,2); \draw[vecteurC] (1,1)--(4,2);
\end{tikzpicture}
\hspace{4pt}
\begin{tikzpicture}[scale=0.42]
  \quadrillage{-0.4}{-0.4}{5.4}{4.4}
  \pt{3,4}{above}{T}
  \draw[vecteur] (0,4)--(2,3); \draw[vecteurB] (2,2)--(1,0);
  \draw[aide] (3,4)--(5,3)--(4,1); \draw[vecteurC] (3,4)--(4,1);
\end{tikzpicture}
\end{center}
\rem{Gauche : $\vec{w}$ fait $3$ à droite et $1$ vers le haut ; droite : $1$ à droite et
$3$ vers le bas.}
\end{exo}

\begin{exo}[Quand c'est l'origine qui manque]
\begin{center}
\begin{tikzpicture}[scale=0.42]
  \quadrillage{-0.4}{-0.4}{5.4}{4.4}
  \pt{4,2}{below right}{A}
  \draw[vecteur] (0,1)--(1,3); \draw[vecteurB] (2,4)--(4,3);
  \draw[vecteurC] (1,1)--(4,2); \pt{1,1}{below}{N}
\end{tikzpicture}
\hspace{4pt}
\begin{tikzpicture}[scale=0.42]
  \quadrillage{-0.4}{-0.4}{5.4}{4.4}
  \pt{2,2}{below right}{A}
  \draw[vecteur] (0,3)--(2,4); \draw[vecteurB] (4,0)--(5,2);
  \draw[vecteurC] (1,3)--(2,2); \pt{1,3}{above left}{N}
\end{tikzpicture}
\end{center}
\rem{On construit le vecteur à partir de $A$ \emph{en reculant} : $N$ est à l'autre bout.
Gauche : $\vec{u}+\vec{v}$ fait $3$ à droite et $1$ vers le haut, donc $N$ est $3$
carreaux à gauche et $1$ plus bas que $A$. Droite : $\vec{u}-\vec{v}$ fait $1$ à droite
et $1$ vers le bas.}
\end{exo}

\begin{exo}[Avec des longueurs]
1) On reporte $\vect{AC}$ à la suite de $\vect{AB}$, depuis $B$ : on arrive en $D$.\par
2) $\vect{BC}-\vect{BA}=\vect{AB}+\vect{BC}$, soit $\vect{AC}$ : on reporte $\vect{AC}$ depuis $B$.\par
3) Par construction, $\vect{BD}=\vect{AC}$ : \res{ABDC \text{ est un parallélogramme}}
(théorème 1) ; ses côtés mesurent $6$~cm et $5$~cm.
\end{exo}

\begin{exo}[Trois vecteurs]
\begin{center}
\begin{tikzpicture}[scale=0.42]
  \quadrillage{-0.4}{-1.4}{8.4}{4.4}
  \pt{1,1}{below left}{A} \pt{5,0}{below}{B}
  \draw[vecteur] (2,3)--(3,4); \draw[vecteurB] (4,4)--(6,3); \draw[vecteurC] (8,1)--(7,3);
  \draw[aide] (1,1)--(2,2)--(4,1); \draw[vecteur] (1,1)--(4,1);
  \draw[aide] (2,2)--(1,4); \draw[vecteurB] (1,1)--(1,4);
  \draw[aide] (5,0)--(7,-1)--(6,1); \draw[vecteurC] (5,0)--(6,1);
\end{tikzpicture}
\end{center}
\rem{$\vec{u}+\vec{v}$ : $3$ à droite ; $\vec{u}+\vec{w}$ : $3$ vers le haut ;
$\vec{v}+\vec{w}$ : $1$ à droite et $1$ vers le haut.}
\end{exo}

\partie{Milieu d'un segment}

\begin{exo}[Une différence, un parallélogramme]
\begin{minipage}[t]{0.5\linewidth}\raggedright
1) $\vect{AC}-\vect{AB}=\vect{BA}+\vect{AC}$, soit $\vect{BC}$. Donc $\vect{AD}=\vect{BC}$ :
\res{ADCB} est un parallélogramme, et ses diagonales $[AC]$ et $[DB]$ ont le même
milieu.\par
2) $\vect{CE}=\vect{CB}+\vect{BA}=\vect{CA}$ : \res{E=A}.
\end{minipage}\hfill
\begin{minipage}[t]{0.46\linewidth}
\vspace{0pt}\centering
\begin{tikzpicture}[scale=0.4]
  \quadrillage{-3.4}{-0.4}{4.4}{3.4}
  \draw[fig] (0,0)--(4,0)--(1,3)--cycle;
  \pt{0,0}{below left}{A} \pt{4,0}{below right}{B} \pt{1,3}{above right}{C}
  \draw[vecteurB] (0,0)--(-3,3); \pt{-3,3}{above}{D}
  \draw[aide] (-3,3)--(1,3) (-3,3)--(4,0) (0,0)--(1,3);
\end{tikzpicture}
\end{minipage}
\end{exo}

\begin{exo}[Trois points à construire]
1) $\vect{AM}=\vect{BC}$ ($ABCM$ est un parallélogramme).\par
2) $\vect{CM}=\vect{CB}+\vect{AB}$ : bout à bout depuis $C$.\par
3) $\vect{AC}+\vect{BA}-\vect{CB}=\vect{BA}+\vect{AC}+\vect{BC}$, soit $\vect{BC}+\vect{BC}$, donc
\res{\vect{BM}=2\vect{BC}}.
\begin{center}
\begin{tikzpicture}[scale=0.36]
  \quadrillage{-3.4}{-3.4}{8.4}{6.4}
  \draw[fig] (0,0)--(4,0)--(1,3)--cycle;
  \pt{0,0}{below left}{A} \pt{4,0}{below right}{B} \pt{1,3}{right}{C}
  \pt{-3,3}{left}{M_1}
  \draw[aide] (1,3)--(4,0)--(8,0); \draw[vecteurB] (1,3)--(8,0); \pt{8,0}{below right}{M_2}
  \draw[vecteurC] (4,0)--(-2,6); \pt{-2,6}{above}{M_3}
\end{tikzpicture}
\end{center}
\rem{Les indices $1$, $2$, $3$ distinguent les points des trois questions.}
\end{exo}

\begin{exo}[Symétrie centrale]
1) $G$ est le milieu de $[EE']$ et de $[FF']$ : les diagonales $[FF']$ et $[EE']$ de
$FEF'E'$ ont le même milieu, donc \res{FEF'E' \text{ est un parallélogramme}}.\par
2) a) Dans ce parallélogramme, $\vect{FE}=\vect{E'F'}$.\par
b) Deux vecteurs égaux ont la même norme : \res{E'F'=EF}.\par
c) La symétrie centrale \res{\text{conserve les longueurs}}.
\end{exo}

\partie{Multiplier un vecteur par un réel}

\begin{exo}[Multiples d'un vecteur]
\begin{center}
\begin{tikzpicture}[scale=0.42]
  \quadrillage{-0.4}{-0.4}{7.4}{5.4}
  \pt{3,3}{left}{O}
  \draw[vecteurB] (3,3)--(6,0) node[pos=0.85,right,font=\scriptsize]{$3\vec{a}$};
  \draw[vecteur] (3,3)--(5,1) node[pos=0.7,left,font=\scriptsize]{$2\vec{a}$};
  \draw[vecteurC] (3,3)--(1,5) node[pos=0.8,right,font=\scriptsize]{$-2\vec{a}$};
  \draw[vecteurC] (3,3)--(2,4) node[pos=0.6,left,font=\scriptsize]{$-\vec{a}$};
\end{tikzpicture}
\end{center}
\rem{$\vec{a}$ fait $1$ à droite et $1$ vers le bas ; $-2\vec{a}$ fait $2$ à gauche et $2$
vers le haut.}
\end{exo}

\begin{exo}[Avec des coefficients]
\begin{center}
\begin{tikzpicture}[scale=0.32]
  \draw[fig] (0,0)--(4,0)--(2,2)--cycle;
  \pt{0,0}{below left}{A} \pt{4,0}{below}{B} \pt{2,2}{left}{C}
  \draw[vecteur] (0,0)--(6,6); \pt{6,6}{above}{M}
  \draw[vecteurB] (2,2)--(3,1); \pt{3,1}{right}{N}
  \draw[vecteurC] (4,0)--(6,0); \pt{6,0}{below right}{P}
\end{tikzpicture}
\end{center}
\rem{$M$ est sur la demi-droite $[AC)$ avec $AM=3AC$ ; $N$ est le milieu de $[CB]$ ;
$P$ est sur la droite $(AB)$, au-delà de $B$, avec $BP=\frac12AB$.}
\end{exo}

\begin{exo}[Sur un quadrillage]
\begin{center}
\begin{tikzpicture}[scale=0.45]
  \quadrillage{-0.4}{-0.4}{8.4}{5.4}
  \pt{0,1}{below left}{A} \pt{4,1}{below}{B} \pt{6,3}{right}{C} \pt{1,4}{above left}{D}
  \draw[vecteur] (0,1)--(4,5) node[pos=0.5,left,font=\scriptsize]{1)};
  \draw[vecteur] (4,1)--(7,2) node[pos=0.6,below,font=\scriptsize]{2)};
  \draw[vecteurB] (1,4)--(2,3) node[pos=0.5,below left,font=\scriptsize]{3)};
  \draw[vecteurB] (6,3)--(5,4) node[pos=0.5,above,font=\scriptsize]{4)};
  \draw[vecteurC] (1,4)--(3,2) node[pos=0.8,right,font=\scriptsize]{5)};
\end{tikzpicture}
\end{center}
\rem{$\vect{BC}$ : $2$ à droite, $2$ vers le haut ; $\vect{AC}$ : $6$ à droite, $2$ vers le
haut ; $\vect{BD}$ : $3$ à gauche, $3$ vers le haut ; $\vect{AB}-\vect{BC}$ : $2$ à
droite, $2$ vers le bas.}
\end{exo}

\partie{Relation de Chasles}

\begin{exo}[Un seul vecteur, si c'est possible]
\deux{a) \res{\vect{BG}}}{b) impossible}
\deux{c) $\vect{KT}+\vect{TU}=\res{\vect{KU}}$}{d) \res{\vect{FS}}}
\deux{e) \res{\vec{0}}}{f) impossible}
\rem{b) Même origine $D$ : Chasles ne s'applique pas. f) $B$ et
$C$ sont différents.}
\end{exo}

\begin{exo}[Simplifier]
\begin{calculs}
\vec{u}&=\vect{BC}+\vect{CA}+\vect{AB}=\vect{BB}\\
\vec{u}&=\res{\vec{0}}\\
\vec{v}&=\vect{CA}+\vect{BC}+\vect{AC}\\
\vec{v}&=\vect{BC}+\vect{CA}+\vect{AC}=\res{\vect{BC}}\\
\vec{w}&=\vect{MC}+\vect{AM}+\vect{CB}\\
\vec{w}&=\vect{AM}+\vect{MC}+\vect{CB}=\res{\vect{AB}}
\end{calculs}
\rem{$\vect{CA}+\vect{AC}=\vect{CC}=\vec{0}$.}
\end{exo}

\begin{exo}[Compléter]
1) $\vect{A\res{C}}$ \quad 2) $\vect{\res{B}A}$ \quad 3) $\vect{A\res{B}}$ \quad
4) $\vect{BC}+\vect{CD}+\vect{DA}=\res{\vect{BA}}$
\end{exo}

\begin{exo}[Peut-on réduire ?]
\deux{a) \res{\vect{MP}}}{b) non (même origine)}
\deux{c) \res{\vect{RS}}}{d) $\vect{SR}+\vect{RO}=\res{\vect{SO}}$}
\deux{e) \res{\vec{0}}}{f) $\vect{AK}+\vect{KB}=\res{\vect{AB}}$}
\end{exo}

\begin{exo}[Une égalité à démontrer]
\begin{calculs}
\vect{OB}-\vect{OA}+\vect{BC}&=\vect{AO}+\vect{OB}+\vect{BC}\\
\vect{OB}-\vect{OA}+\vect{BC}&=\res{\vect{AC}}
\end{calculs}
\rem{$-\vect{OA}=\vect{AO}$, puis on range : $A\to O\to B\to C$.}
\end{exo}

\begin{exo}[Exprimer avec deux vecteurs]
1) a)
\begin{calculs}
\vec{u}&=\vect{CA}+\vect{AB}\\
\vec{u}&=\res{\vect{AB}-\vect{AC}}
\end{calculs}
b)
\begin{calculs}
\vec{u}&=3(\vect{AB}-\vect{AC})-\vect{AB}\\
\vec{u}&=\res{2\vect{AB}-3\vect{AC}}
\end{calculs}
2) a)
\begin{calculs}
\vec{v}&=\vect{BC}+\vect{CA}+\vect{CA}\\
\vec{v}&=\res{\vect{BC}+2\vect{CA}}
\end{calculs}
b) $\vect{AB}=\vect{AC}+\vect{CB}=-\vect{CA}-\vect{BC}$ et $\vect{AC}=-\vect{CA}$ :
\begin{calculs}
\vec{v}&=-\vect{CA}-\vect{BC}+2\vect{BC}-\vect{CA}\\
\vec{v}&=\res{\vect{BC}-2\vect{CA}}
\end{calculs}
\end{exo}

\begin{exo}[Différences sur un quadrillage]
\begin{center}
\begin{tikzpicture}[scale=0.42]
  \quadrillage{-0.4}{-2.4}{8.4}{4.4}
  \pt{0,1}{left}{A} \pt{4,0}{below}{B} \pt{5,3}{above right}{C} \pt{1,3}{above}{D}
  \draw[vecteur] (0,1)--(1,4) node[pos=0.6,left,font=\scriptsize]{a)};
  \draw[vecteurB] (0,1)--(4,0) node[pos=0.5,above,font=\scriptsize]{b)};
  \draw[vecteurC] (0,1)--(3,-2) node[pos=0.5,left,font=\scriptsize]{c)};
  \draw[vecteur] (0,1)--(8,0) node[pos=0.85,below,font=\scriptsize]{d)};
\end{tikzpicture}
\end{center}
\rem{a) $\vect{BC}$ ; b) $\vect{AD}+\vect{DB}=\vect{AB}$ : on arrive en $B$ ;
c) $\vect{DB}$ ; d) $\vect{AC}+\vect{DB}$ : $8$ à droite et $1$ vers le bas.}
\end{exo}

\end{document}
